\documentclass[10pt,a4paper]{amsart}
\usepackage{amsmath,amssymb,amsthm,graphicx}
% Retained minimal preamble for the Stacks Project chapter flat.tex (More on Flatness).
% Verbatim excerpt of the chapter preamble: Xy-pic setup (native preamble lines 14-18) and the native
% theorem declarations with the mathematical macros (native preamble lines 154-217). Copied mechanically;
% nothing added, removed or redefined. The wrapper class and the amsmath/amssymb/amsthm/graphicx base
% packages are supplied by the assembler, so the class/externaldocument part of the native preamble is omitted.
% For commutative diagrams we use Xy-pic
\usepackage[all]{xy}

% We use 2cell for 2-commutative diagrams.
\xyoption{2cell}
% Theorem environments.
%
\theoremstyle{plain}
\newtheorem{theorem}[subsection]{Theorem}
\newtheorem{proposition}[subsection]{Proposition}
\newtheorem{lemma}[subsection]{Lemma}

\theoremstyle{definition}
\newtheorem{definition}[subsection]{Definition}
\newtheorem{example}[subsection]{Example}
\newtheorem{exercise}[subsection]{Exercise}
\newtheorem{situation}[subsection]{Situation}

\theoremstyle{remark}
\newtheorem{remark}[subsection]{Remark}
\newtheorem{remarks}[subsection]{Remarks}

\numberwithin{equation}{subsection}

% Macros
%
\def\lim{\mathop{\mathrm{lim}}\nolimits}
\def\colim{\mathop{\mathrm{colim}}\nolimits}
\def\Spec{\mathop{\mathrm{Spec}}}
\def\Hom{\mathop{\mathrm{Hom}}\nolimits}
\def\Ext{\mathop{\mathrm{Ext}}\nolimits}
\def\SheafHom{\mathop{\mathcal{H}\!\mathit{om}}\nolimits}
\def\SheafExt{\mathop{\mathcal{E}\!\mathit{xt}}\nolimits}
\def\Sch{\mathit{Sch}}
\def\Mor{\mathop{\mathrm{Mor}}\nolimits}
\def\Ob{\mathop{\mathrm{Ob}}\nolimits}
\def\Sh{\mathop{\mathit{Sh}}\nolimits}
\def\NL{\mathop{N\!L}\nolimits}
\def\CH{\mathop{\mathrm{CH}}\nolimits}
\def\proetale{{pro\text{-}\acute{e}tale}}
\def\etale{{\acute{e}tale}}
\def\QCoh{\mathit{QCoh}}
\def\Ker{\mathop{\mathrm{Ker}}}
\def\Im{\mathop{\mathrm{Im}}}
\def\Coker{\mathop{\mathrm{Coker}}}
\def\Coim{\mathop{\mathrm{Coim}}}

% Boxtimes
%
\DeclareMathSymbol{\boxtimes}{\mathbin}{AMSa}{"02}

%
% Macros for moduli stacks/spaces
%
\def\QCohstack{\mathcal{QC}\!\mathit{oh}}
\def\Cohstack{\mathcal{C}\!\mathit{oh}}
\def\Spacesstack{\mathcal{S}\!\mathit{paces}}
\def\Quotfunctor{\mathrm{Quot}}
\def\Hilbfunctor{\mathrm{Hilb}}
\def\Curvesstack{\mathcal{C}\!\mathit{urves}}
\def\Polarizedstack{\mathcal{P}\!\mathit{olarized}}
\def\Complexesstack{\mathcal{C}\!\mathit{omplexes}}
% \Pic is the operator that assigns to X its picard group, usage \Pic(X)
% \Picardstack_{X/B} denotes the Picard stack of X over B
% \Picardfunctor_{X/B} denotes the Picard functor of X over B
\def\Pic{\mathop{\mathrm{Pic}}\nolimits}
\def\Picardstack{\mathcal{P}\!\mathit{ic}}
\def\Picardfunctor{\mathrm{Pic}}
\def\Deformationcategory{\mathcal{D}\!\mathit{ef}}

% --- appended 2026-10-12 (Chief): the project own definitions of the reference/slogan/history
% display environments, copied mechanically from the shared Stacks preamble lines 8-12
% (\usepackage{verbatim} plus the three \newenvironment lines, which the chapter preamble keeps
% verbatim from the project preamble). Nothing added, removed or redefined: these environments
% carry no mathematical content and the project itself typesets them as comments. Without these
% three lines a counted statement whose own text uses them fails with Environment slogan undefined.
\usepackage{verbatim}
\newenvironment{reference}{\comment}{\endcomment}
\newenvironment{slogan}{\comment}{\endcomment}
\newenvironment{history}{\comment}{\endcomment}
\usepackage{cleveref}
\crefname{originallemma}{original lemma}{original lemmas}
\Crefname{originallemma}{Original lemma}{Original lemmas}
\makeatletter
\newlabel{algebra-lemma-countgen-projective}{{059X}{source}{}{sourcepointer-1}{}}
\newlabel{algebra-lemma-countgen-projective@cref}{{[originallemma][0][]059X}{[1][0][]source}{}{}{}}
\makeatother
\crefname{originallemma}{original lemma}{original lemmas}
\Crefname{originallemma}{Original lemma}{Original lemmas}
\makeatletter
\newlabel{algebra-lemma-pure-submodule-ML}{{059N}{source}{}{sourcepointer-2}{}}
\newlabel{algebra-lemma-pure-submodule-ML@cref}{{[originallemma][0][]059N}{[1][0][]source}{}{}{}}
\makeatother
\crefname{originallemma}{original lemma}{original lemmas}
\Crefname{originallemma}{Original lemma}{Original lemmas}
\makeatletter
\newlabel{descent-lemma-locally-projective-descends}{{05JZ}{source}{}{sourcepointer-3}{}}
\newlabel{descent-lemma-locally-projective-descends@cref}{{[originallemma][0][]05JZ}{[1][0][]source}{}{}{}}
\makeatother
\crefname{originallemma}{original lemma}{original lemmas}
\Crefname{originallemma}{Original lemma}{Original lemmas}
\makeatletter
\newlabel{lemma-affine-locally-projective-pure}{{05K6}{source}{}{sourcepointer-4}{}}
\newlabel{lemma-affine-locally-projective-pure@cref}{{[originallemma][0][]05K6}{[1][0][]source}{}{}{}}
\makeatother
\crefname{originallemma}{original lemma}{original lemmas}
\Crefname{originallemma}{Original lemma}{Original lemmas}
\makeatletter
\newlabel{lemma-base-change-universally-flat}{{05DH}{source}{}{sourcepointer-5}{}}
\newlabel{lemma-base-change-universally-flat@cref}{{[originallemma][0][]05DH}{[1][0][]source}{}{}{}}
\makeatother
\crefname{originallemma}{original lemma}{original lemmas}
\Crefname{originallemma}{Original lemma}{Original lemmas}
\makeatletter
\newlabel{lemma-criterion}{{05L4}{source}{}{sourcepointer-6}{}}
\newlabel{lemma-criterion@cref}{{[originallemma][0][]05L4}{[1][0][]source}{}{}{}}
\makeatother
\crefname{originallemma}{original lemma}{original lemmas}
\Crefname{originallemma}{Original lemma}{Original lemmas}
\makeatletter
\newlabel{lemma-finite-presentation-flat-along-fibre}{{05KW}{source}{}{sourcepointer-7}{}}
\newlabel{lemma-finite-presentation-flat-along-fibre@cref}{{[originallemma][0][]05KW}{[1][0][]source}{}{}{}}
\makeatother
\crefname{originallemma}{original lemma}{original lemmas}
\Crefname{originallemma}{Original lemma}{Original lemmas}
\makeatletter
\newlabel{lemma-finite-type-flat-pure-along-fibre-is-universal}{{05L5}{source}{}{sourcepointer-8}{}}
\newlabel{lemma-finite-type-flat-pure-along-fibre-is-universal@cref}{{[originallemma][0][]05L5}{[1][0][]source}{}{}{}}
\makeatother
\crefname{originallemma}{original lemma}{original lemmas}
\Crefname{originallemma}{Original lemma}{Original lemmas}
\makeatletter
\newlabel{more-morphisms-lemma-relative-assassin-constructible}{{05KR}{source}{}{sourcepointer-9}{}}
\newlabel{more-morphisms-lemma-relative-assassin-constructible@cref}{{[originallemma][0][]05KR}{[1][0][]source}{}{}{}}
\makeatother
\crefname{originallemma}{original lemma}{original lemmas}
\Crefname{originallemma}{Original lemma}{Original lemmas}
\makeatletter
\newlabel{morphisms-lemma-constructible-containing-open}{{05LW}{source}{}{sourcepointer-10}{}}
\newlabel{morphisms-lemma-constructible-containing-open@cref}{{[originallemma][0][]05LW}{[1][0][]source}{}{}{}}
\makeatother
\title{Projectivity of the pushforward of a flat finitely presented module is equivalent to purity relative to the base}
\author{The Stacks Project Authors}
\date{stacks-project:flat.tex:local-snapshot-20261011; GNU FDL 1.2 or later}
\begin{document}
\maketitle
\noindent\textit{Editorial note.} All delivered mathematical text is the author's own, copied from the pinned Stacks Project chapter file flat.tex; no statement or proof step was rewritten, reordered or removed.\par\medskip
\noindent\textit{Editorial note.} Preamble provenance: the chapter's own retained minimal preamble (sources/stacks-flat/retained-preamble.tex, lines 1-85) is a verbatim excerpt of the native Xy-pic setup and of the native theorem declarations and mathematical macros, followed by the project's own reference/slogan/history display-environment definitions copied mechanically from the shared preamble; the wrapper class and the amsmath/amssymb/amsthm/graphicx base packages are supplied by the assembler, and extra\_packages is left empty because everything needed is already in the retained preamble. No macro was invented or redefined.\par\medskip
\noindent\textit{Editorial note.} Same-chapter pointers are located at their \textbackslash{}label lines inside the pinned chapter file; cross-chapter pointers are located in the project's own tag list staged in the snapshot (sources/stacks-chow/tags.txt), at the verbatim lines carrying their canonical tags.\par\medskip
\noindent\textit{Allowed standard prerequisites (editorial).}\begin{itemize}
\item Purity of a module with locally projective global sections. Same chapter flat.tex, label lemma-affine-locally-projective-pure, invoked at proof lines 5599 and 5622. Gives the easy implication (1) implies (2), and shows that the module pulled back along the etale neighbourhood is universally pure over the base.
\item Universal purity from purity along the base, and fpqc descent of local projectivity. Same chapter flat.tex, label lemma-finite-type-flat-pure-along-fibre-is-universal (proof line 5602); Descent, label descent-lemma-locally-projective-descends (proof line 5604). Makes purity stable under base change, which is what allows the etale local computation, and reduces projectivity of the pushforward to fpqc local projectivity.
\item The etale neighbourhood with projective global sections, the purity criterion, constructibility of the relative assassin, and the criterion for universal injectivity. Same chapter flat.tex, labels lemma-finite-presentation-flat-along-fibre (proof line 5609), lemma-criterion (proof line 5624) and lemma-base-change-universally-flat (proof line 5641); More on Morphisms, label more-morphisms-lemma-relative-assassin-constructible (proof line 5633); Morphisms, label morphisms-lemma-constructible-containing-open (proof line 5636). Produces the affine etale neighbourhood, places the relative associated points inside the image of the etale morphism after shrinking, and converts that inclusion into universal injectivity of the map of global sections.
\item The Mittag-Leffler criterion for purity of a submodule and projectivity of a countably generated flat module. Algebra, labels algebra-lemma-pure-submodule-ML (proof line 5647) and algebra-lemma-countgen-projective (proof line 5653). Turns universal injectivity of the map of global sections into the Mittag-Leffler property, which combined with countable generation and flatness gives projectivity of the pushforward.
\end{itemize}\medskip
\ifdefined\hypertarget\hypertarget{sourcepointer-1}{}\fi
\noindent\textit{Source reference.} Original lemma 059X: sources/stacks-chow/tags.txt, source lines 6830--6830. Navigation pointer for \textbackslash{}ref\{algebra-lemma-countgen-projective\} used inside the counted span (line 5653); the argument that consumes it is retained in the span.\par\medskip
\ifdefined\hypertarget\hypertarget{sourcepointer-2}{}\fi
\noindent\textit{Source reference.} Original lemma 059N: sources/stacks-chow/tags.txt, source lines 6821--6821. Navigation pointer for \textbackslash{}ref\{algebra-lemma-pure-submodule-ML\} used inside the counted span (line 5647); the argument that consumes it is retained in the span.\par\medskip
\ifdefined\hypertarget\hypertarget{sourcepointer-3}{}\fi
\noindent\textit{Source reference.} Original lemma 05JZ: sources/stacks-chow/tags.txt, source lines 7183--7183. Navigation pointer for \textbackslash{}ref\{descent-lemma-locally-projective-descends\} used inside the counted span (line 5604); the argument that consumes it is retained in the span.\par\medskip
\ifdefined\hypertarget\hypertarget{sourcepointer-4}{}\fi
\noindent\textit{Source reference.} Original lemma 05K6: sources/stacks-flat/pinned-original-source.tex, source lines 5148--5148. Navigation pointer for \textbackslash{}ref\{lemma-affine-locally-projective-pure\} used inside the counted span (line 5599); the argument that consumes it is retained in the span.\par\medskip
\ifdefined\hypertarget\hypertarget{sourcepointer-5}{}\fi
\noindent\textit{Source reference.} Original lemma 05DH: sources/stacks-flat/pinned-original-source.tex, source lines 1907--1907. Navigation pointer for \textbackslash{}ref\{lemma-base-change-universally-flat\} used inside the counted span (line 5641); the argument that consumes it is retained in the span.\par\medskip
\ifdefined\hypertarget\hypertarget{sourcepointer-6}{}\fi
\noindent\textit{Source reference.} Original lemma 05L4: sources/stacks-flat/pinned-original-source.tex, source lines 5197--5197. Navigation pointer for \textbackslash{}ref\{lemma-criterion\} used inside the counted span (line 5624); the argument that consumes it is retained in the span.\par\medskip
\ifdefined\hypertarget\hypertarget{sourcepointer-7}{}\fi
\noindent\textit{Source reference.} Original lemma 05KW: sources/stacks-flat/pinned-original-source.tex, source lines 3342--3342. Navigation pointer for \textbackslash{}ref\{lemma-finite-presentation-flat-along-fibre\} used inside the counted span (line 5609); the argument that consumes it is retained in the span.\par\medskip
\ifdefined\hypertarget\hypertarget{sourcepointer-8}{}\fi
\noindent\textit{Source reference.} Original lemma 05L5: sources/stacks-flat/pinned-original-source.tex, source lines 5276--5276. Navigation pointer for \textbackslash{}ref\{lemma-finite-type-flat-pure-along-fibre-is-universal\} used inside the counted span (line 5602); the argument that consumes it is retained in the span.\par\medskip
\ifdefined\hypertarget\hypertarget{sourcepointer-9}{}\fi
\noindent\textit{Source reference.} Original lemma 05KR: sources/stacks-chow/tags.txt, source lines 7210--7210. Navigation pointer for \textbackslash{}ref\{more-morphisms-lemma-relative-assassin-constructible\} used inside the counted span (line 5633); the argument that consumes it is retained in the span.\par\medskip
\ifdefined\hypertarget\hypertarget{sourcepointer-10}{}\fi
\noindent\textit{Source reference.} Original lemma 05LW: sources/stacks-chow/tags.txt, source lines 7250--7250. Navigation pointer for \textbackslash{}ref\{morphisms-lemma-constructible-containing-open\} used inside the counted span (line 5636); the argument that consumes it is retained in the span.\par\medskip
\begin{proposition}
\label{proposition-finite-presentation-flat-pure-is-projective}
Let $f : X \to S$ be an affine, finitely presented morphism of schemes.
Let $\mathcal{F}$ be a quasi-coherent $\mathcal{O}_X$-module of
finite presentation, flat over $S$. Then the following
are equivalent
\begin{enumerate}
\item $f_*\mathcal{F}$ is locally projective on $S$, and
\item $\mathcal{F}$ is pure relative to $S$.
\end{enumerate}
In particular, given a ring map $A \to B$ of finite presentation and
a finitely presented $B$-module $N$ flat over $A$ we have:
$N$ is projective as an $A$-module if and only if $\widetilde{N}$
on $\Spec(B)$ is pure relative to $\Spec(A)$.
\end{proposition}

\begin{proof}
The implication (1) $\Rightarrow$ (2) is
Lemma \ref{lemma-affine-locally-projective-pure}.
Assume $\mathcal{F}$ is pure relative to $S$.
Note that by
Lemma \ref{lemma-finite-type-flat-pure-along-fibre-is-universal}
this implies $\mathcal{F}$ remains pure after any base change. By
Descent, Lemma \ref{descent-lemma-locally-projective-descends}
it suffices to prove $f_*\mathcal{F}$ is fpqc locally projective on $S$.
Pick $s \in S$. We will prove that the restriction of
$f_*\mathcal{F}$ to an \'etale neighbourhood of $s$ is locally projective.
Namely, by
Lemma \ref{lemma-finite-presentation-flat-along-fibre},
after replacing $S$ by an affine elementary \'etale
neighbourhood of $s$, we may assume there exists a diagram
$$
\xymatrix{
X \ar[dr] & & X' \ar[ll]^g \ar[ld] \\
& S &
}
$$
of schemes affine and of finite presentation over $S$,
where $g$ is \'etale, $X_s \subset g(X')$, and with
$\Gamma(X', g^*\mathcal{F})$ a projective $\Gamma(S, \mathcal{O}_S)$-module.
Note that in this case $g^*\mathcal{F}$ is universally pure over $S$, see
Lemma \ref{lemma-affine-locally-projective-pure}.
Hence by
Lemma \ref{lemma-criterion}
we see that the open $g(X')$ contains the points of
$\text{Ass}_{X/S}(\mathcal{F})$ lying over $\Spec(\mathcal{O}_{S, s})$.
Set
$$
E = \{t \in S \mid \text{Ass}_{X_t}(\mathcal{F}_t) \subset g(X') \}.
$$
By
More on Morphisms,
Lemma \ref{more-morphisms-lemma-relative-assassin-constructible}
$E$ is a constructible subset of $S$. We have seen that
$\Spec(\mathcal{O}_{S, s}) \subset E$. By
Morphisms, Lemma \ref{morphisms-lemma-constructible-containing-open}
we see that $E$ contains an open neighbourhood of $s$. Hence after
replacing $S$ by an affine neighbourhood of $s$ we may assume that
$\text{Ass}_{X/S}(\mathcal{F}) \subset g(X')$.
By
Lemma \ref{lemma-base-change-universally-flat}
this means that
$$
\Gamma(X, \mathcal{F}) \longrightarrow \Gamma(X', g^*\mathcal{F})
$$
is $\Gamma(S, \mathcal{O}_S)$-universally injective.
By Algebra, Lemma \ref{algebra-lemma-pure-submodule-ML}
we conclude that $\Gamma(X, \mathcal{F})$ is Mittag-Leffler as an
$\Gamma(S, \mathcal{O}_S)$-module. Since
$\Gamma(X, \mathcal{F})$ is countably generated and flat as a
$\Gamma(S, \mathcal{O}_S)$-module, we conclude
it is projective by
Algebra, Lemma \ref{algebra-lemma-countgen-projective}.
\end{proof}

\end{document}
